\documentclass[11pt]{article}

% ----------------------------
% Packages
% ----------------------------
\usepackage[margin=1in]{geometry}
\usepackage{amsmath, amssymb, amsthm}
\usepackage{graphicx}
\usepackage{hyperref}
\usepackage{enumitem}
\usepackage{mathtools}
\usepackage{booktabs}
\usepackage{float}
\usepackage{authblk}

% ----------------------------
% Theorem Environments
% ----------------------------
\theoremstyle{definition}
\newtheorem{definition}{Definition}[section]

\theoremstyle{plain}
\newtheorem{theorem}{Theorem}[section]
\newtheorem{lemma}{Lemma}[section]

\theoremstyle{remark}
\newtheorem{remark}{Remark}[section]

% ----------------------------
% Title Info
% ----------------------------
\title{\textbf{Pricing Multi-Dimension Bermudan Options LSM-FNN vs LSM-Poly}}

\author[1]{Brennan Davenport}
\author[1]{Nathaniel Levy}

\affil[1]{University of Georgia, USA \\
\texttt{\{Brennan.Davenport@uga.edu, Nathaniel.Levy@uga.edu\}}}

\date{\today}


% ----------------------------
% Document
% ----------------------------
\begin{document}

\maketitle

\begin{abstract}
We study \ldots  
Briefly state the problem, your approach, and the main results.
\end{abstract}

\section{Introduction}
\label{sec:intro}

Explain the problem, why it matters, and what your contribution is.
A common structure:
\begin{itemize}
    \item What problem are you solving?
    \item Why is it hard / interesting?
    \item What do you contribute?
\end{itemize}

\section{Related Work}
\label{sec:related}

Discuss prior approaches and how your work differs.
Cite papers using \cite{example_reference}.

\section{Mathematical Setup}
\label{sec:setup}


Let us first introduce a few definitions and assumptions. Our goal is to price multiple Bermudan options. 
For concreteness, we will refer to this collection as a basket of Bermudan options, or simply a basket option.

To make the model well-defined, we assume that the risk-free interest rate satisfies \(r \ge 0\). 
Next, we define a set of exercise times
\[
\mathcal{T} \subset [0, T],
\] 
which represents the dates at which the Bermudan option may be exercised.

We now define the payoff function. Let  
\[
h(S_t) : \mathbb{R}_+^d \to \mathbb{R}
\]  
denote the payoff function, where \(S_t \in \mathbb{R}_+^d\) is the vector of underlying asset prices 
at time \(t\). Furthermore, let \(d \in \mathbb{Z}_{\geq 1}\) be the number dimension of this basket.
The form of \(h\) depends on whether the option is a put or a call. In this project, we 
assume that the payoff is based on the average of the underlying assets.

\begin{itemize}
    \item PUT option:
    \[
    h(S_t) := \max\Big(K - \frac{1}{d} \sum_{j=1}^d S_t^{(j)}, 0\Big)
    \]
    \item CALL option:
    \[
    h(S_t) := \max\Big(\frac{1}{d} \sum_{j=1}^d S_t^{(j)} - K, 0\Big)
    \]
\end{itemize}

Here, \(K \in \mathbb{R}_+\) is the strike price.

With the payoff function defined, we can now define the value function for the basket option. Let
\[
V(t, S_t) : \mathcal{T} \times \mathbb{R}_+^d \to \mathbb{R}
\]  
denote the value function of a \(d\)-dimensional Bermudan option. 
The vector \(S_t\) represents the asset prices at time \(t\).

Now that we have defined the value function \(V\), we describe its form using the Bellman equation:
\[
V(t_i, S_{t_i}) = \max \Big( h(S_{t_i}),\; e^{-r (t_{i+1} - t_i)} \mathbb{E}[ V(t_{i+1}, S_{t_{i+1}}) \mid S_{t_i} ] \Big),
\]

where \(t_i \in \mathcal{T}\). This provides an inductive definition of \(V\). And we know that at the final 
exercise date \(m = \max(\mathcal{T})\), we have
\[
V(m, S_m) = h(S_m),
\]
since there is no continuation beyond the last exercise date.

We now focus on the second term of the Bellman equation:
\[
C_i(S_{t_i}) := e^{-r (t_{i+1} - t_i)} \mathbb{E}[ V(t_{i+1}, S_{t_{i+1}}) \mid S_{t_i} ].
\]
The factor \(e^{-r (t_{i+1}-t_i)}\) represents the discounted cash flow, which is straightforward
(assuming a risk-free model). The challenging part is the conditional expectation
\[
\mathbb{E}[ V(t_{i+1}, S_{t_{i+1}}) \mid S_{t_i} ],
\]
which we refer to as the continuation value. In high dimensions, it is infeasible to compute this exactly using a lookup table.


In the classical Longstaff-Schwartz method, this conditional expectation is approximated using polynomial regression:
\[
\mathbb{E}[ V(t_{i+1}, S_{t_{i+1}}) \mid S_{t_i} ] \approx \sum_{\alpha} \beta_\alpha \, \phi_\alpha(S_{t_i}),
\]
where \(\{\phi_\alpha\}\) is a set of polynomial basis functions of the 
current asset prices, and \(\beta_\alpha\) are the regression coefficients.


The goal of this research is to replace the polynomial regression approximation with a feedforward neural network (FNN):
\[
\mathbb{E}[ V(t_{i+1}, S_{t_{i+1}}) \mid S_{t_i} ] \approx \text{FNN}(S_{t_i}; \theta),
\]
where \(\theta\) represents the trainable parameters of the neural network. 
This approach allows us to efficiently approximate the continuation value even in high-dimensional settings.



\section{Methodology}
\label{sec:method}

Describe your algorithm / model.

\subsection{Model}
Explain the model intuition first, then equations.

\subsection{Algorithm}
Pseudo-code or step-by-step explanation.

\section{Experiments / Results}
\label{sec:results}

\subsection{Experimental Setup}
Data, parameters, assumptions.

\subsection{Results}
Tables, plots, and discussion.

\begin{table}[H]
\centering
\begin{tabular}{ccc}
\toprule
Method & Price & Runtime \\
\midrule
LSM (Poly) & 10.42 & 0.31s \\
LSM (FNN)  & 10.57 & 0.48s \\
\bottomrule
\end{tabular}
\caption{Pricing results}
\end{table}

\section{Discussion}
\label{sec:discussion}

Interpret results, tradeoffs, limitations.

\section{Conclusion}
\label{sec:conclusion}

Summarize findings and suggest future work.

% ----------------------------
% Bibliography
% ----------------------------
\bibliographystyle{plain}
\bibliography{references}

\end{document}
